\section{Closed-Return Correction}

\subsection{Closed-return structure}

The first-order bridge equation is
\[
x=B+\frac{K}{x}.
\]
Its positive root is
\[
x_{\rm quad}=137.03600027236\ldots.
\]
This first-order result is structurally constrained but not sufficient at
modern comparison precision. The admitted cover-mediated feedback has a
first closed return on the scalar two-node bridge graph. That return
supplies the second-order correction
\[
\boxed{
\Delta x_{\rm 2nd}
=
-\frac{K N_m}{(N_e-1)x^3}.
}
\]

\subsection{The two-node feedback graph}

The relevant feedback graph contains two transport nodes:
\[
\text{EM bridge}
\]
and
\[
\text{paired magnetic source}.
\]
The cover map
\[
\pi:Z_8\to Z_4
\]
is channel-internal, not an independent third transport node. A structure
would qualify as an independent transport node only if it had independent
state space, independent transport capacity, sequential multiplicative
factorization, and independent normalization.

The cover map \(\pi\) fails these criteria. It is a morphism, not a
carrier object with its own independent position space. Its kernel is
internal to \(Z_8\), and elements in a fiber are co-accessed by
base-level operations factoring through \(\pi\). Additive cover-to-base
readout therefore has the signature of mechanism-distinct routes through
a channel, not sequential traversal through a separate transport node.

Thus the active scalar feedback graph is
\[
\text{EM bridge}
\longleftrightarrow
\text{paired magnetic source}.
\]

\subsection{First closed return}

The first-order term
\[
+\frac{K}{x}
\]
is an outward cover-mediated feedback contribution from the paired
magnetic source into the electromagnetic bridge. The first closed return
is the shortest circuit that leaves the bridge and returns to it through
the paired magnetic source:
\[
\text{EM bridge}
\rightarrow
\text{paired magnetic source}
\rightarrow
\text{EM bridge}.
\]
This return circuit requires two additional bridge traversals beyond the
first-order feedback term. Each traversal is diluted over the effective
bridge capacity \(x\), so the two additional traversals contribute
\[
\frac1{x^2}.
\]
The return therefore has the scale
\[
\frac{K}{x}\cdot\frac1{x^2}
=
\frac{K}{x^3}.
\]

\subsection{Return coefficient}

The return coefficient is
\[
\frac{N_m}{N_e-1}.
\]
The numerator \(N_m\) is the magnetic reception capacity of the paired
magnetic side. The denominator \(N_e-1\) is the electric-cover complement
capacity remaining after the first-order source activation.

This denominator follows from co-access plus transport accounting. A
base-level activation through
\[
\pi:Z_8\to Z_4
\]
co-accesses the full fiber
\[
\pi^{-1}(m),
\]
but it does so as one carrier-level transport event, not as independent
electric activations for each element of the fiber. Therefore the
first-order event removes one transport-event unit from the electric-cover
transport capacity, leaving
\[
N_e-1.
\]
It does not remove \(d\) independent electric sites, so the denominator
is not \(N_e-d\). It is also not \(N_e\), because the first-order source
activation has already committed one electric-cover transport event
before the return is formed.

For the compact fiber,
\[
N_m=4,
\qquad
N_e=8,
\]
and therefore
\[
\frac{N_m}{N_e-1}
=
\frac47.
\]

\subsection{Negative sign}

The sign of the closed return is fixed by its role as derivative
self-response of the first-order feedback term. The first-order feedback
contribution is
\[
g(x)=\frac{K}{x}.
\]
Since
\[
K>0,
\qquad
x>0,
\]
we have
\[
g'(x)=-\frac{K}{x^2}<0.
\]
Thus, if the closed return is the derivative self-response of the
first-order feedback to the bridge capacity that feedback modifies, it
must reduce the outward first-order contribution. The path factors
\[
\frac1{x^2}
\]
and
\[
\frac{N_m}{N_e-1}
\]
are positive. They affect the magnitude but not the sign. Therefore
\[
\Delta x_{\rm 2nd}<0.
\]

\subsection{Second-order correction}

Combining the first-order admitted feedback, the two additional bridge
traversals, the return coefficient, and the negative derivative
self-response sign gives
\[
\Delta x_{\rm 2nd}
=
-\frac{K}{x}
\cdot
\frac1{x^2}
\cdot
\frac{N_m}{N_e-1}.
\]
Therefore
\[
\boxed{
\Delta x_{\rm 2nd}
=
-\frac{K N_m}{(N_e-1)x^3}.
}
\]

For the compact-fiber values,
\[
K=\frac{74}{15},
\qquad
N_m=4,
\qquad
N_e=8,
\]
the coefficient is
\[
\frac{K N_m}{N_e-1}
=
\frac{(74/15)\cdot4}{7}
=
\frac{296}{105}.
\]
Thus the second-order correction is
\[
\boxed{
\Delta x_{\rm 2nd}
=
-\frac{296}{105x^3}.
}
\]

The measured fine-structure constant is not used to set
\[
\frac47
\]
or any other coefficient in the corrected fixed-point equation.

\subsection{Corrected fixed point}

Adding the closed-return correction to the first-order equation gives
\[
x
=
B+\frac{K}{x}
-
\frac{K N_m}{(N_e-1)x^3}.
\]
Substituting
\[
B=137,
\qquad
K=\frac{74}{15},
\qquad
N_m=4,
\qquad
N_e=8,
\]
we obtain
\[
x
=
137+\frac{74}{15x}
-
\frac{296}{105x^3}.
\]
This is the second-order compact-fiber bridge equation.

\subsection{Status of the correction}

The admitted second-order correction is conditionally unique within the
current scalar two-node bridge graph. At order \(O(x^{-3})\), possible
sources would include a second independent first-order feedback source,
a new carrier pool, a direct cubic mechanism, base-feedback interference,
a non-scalar correction, or higher-order leakage. Under the current scalar
two-node graph, none of these alternatives is admitted: the closed return
of the already-admitted first-order feedback is the only same-order scalar
closed-return correction.

This is a conditional uniqueness claim. It is relative to the current
scalar bridge graph and the stated transport-node, source/complement,
derivative self-response, and scalar-readout rules. Future radiation-side
feedback mechanisms would require re-audit, but no such mechanism is
admitted in the present scalar bridge graph.

If repeated closed returns are admitted, each additional round trip is
suppressed by at least
\[
\frac{C}{x^2},
\qquad
C=\frac{N_m}{N_e-1}=\frac47.
\]
For
\[
x\approx137.036,
\]
this suppression factor is approximately
\[
3.0\times10^{-5}.
\]
Thus higher-order repeated returns, if present, are operationally small
at the current comparison scale. A complete all-orders return theorem
remains open; the present paper uses only the first closed-return
correction.